For the validated TD MVP SP1 relation, soundness reduces to two assumptions:
the SP1 proof system soundly proves execution of the guest program, and
SHA-256 commitments are collision resistant for the artifact sizes considered.
Under these assumptions, an accepting proof implies that the guest relation
accepted the public input and private witness.

For committed TD, an accepting proof binds the transition to the public dataset
root and binds the public claimed target/loss to deterministic fixed-point
recomputation. A prover cannot change the transition, reward, done flag, leaf
index, Merkle path, target witness, TD error, or loss without causing rejection,
except with a hash collision or proof-system failure.

For distinct minibatch TD, the Python relation check extends itemwise and also
covers batch structure. The public ordered index list must match the private
items, duplicate indices are rejected, and the claimed batch loss must equal
the integer average of recomputed per-item losses. This blocks attacks that
replace one item, shuffle items while preserving a multiset, or adjust the
claimed average in the Python oracle path.

For model-grounded forward-TD Python cases, the selected Q-values are no longer
free TD witnesses. They are anchored to committed quantized weights and checked
through MLP forward execution, ReLU masks, online argmax semantics, and
target-network value selection.

For the tiny one-step SGD Python case, the check additionally binds the
claimed post-update model to the checked pre-update model, learning rate,
SmoothL1 derivative branch, gradient tensors, and fixed-point delta tensors.
This is update evidence for one small step, not a proof that repeated optimizer
steps, target-network synchronizations, or sampling decisions formed a complete
training run.

The guarantees are local to the checked relations. They do not authenticate the
data-generation process before the Merkle root was chosen, and they do not prove
that unchecked training iterations or optimizer steps happened.
